\subsubsection{Anti-sense and Intergenic Transcription}

Each PacBio isoform cluster was labelled base-by-base against the canonical~\syn~gene model, and its sense, antisense, and intergenic coverage fractions were computed.
Antisense coverage was carried by 2\% of the isoforms (345 of those with $\geq 10$ supporting reads), which collapsed to 90 representative clusters after merging near-identical ends within 10~bp.
These were classified by the arrangement of their sense and antisense segments into three cases: antisense-initiated transcripts from a spurious promoter in the AT-rich genome (anti $\rightarrow \ldots$), sense transcripts reading through into an antisense gene at an operon boundary (sense $\rightarrow$ anti), and read-through transcripts with an antisense gene embedded between sense genes (sense $\rightarrow$ anti $\rightarrow$ sense).
Spurious-promoter transcripts accounted for 60 of the 90 clusters (67\%), and read-through transcripts, including 4 embedded cases, for the remaining 30 (33\%).

\subsubsection{Novel Open-Reading Frames from Full-length RNA Isoforms}

The abnormal isoforms, defined as those with $>$10 supporting reads and an abnormal fraction ($f_\mathrm{antisense} + f_\mathrm{intergenic}$) $>$ 0.2, gave 816 isoforms as input to the ORF-prediction pipeline.

All candidate start codons within the 816 isoforms were evaluated using OSTIR~\cite{roots_ostir_2021}, which predicts translation initiation rates (TIR) by modelling Shine--Dalgarno (SD) complementarity to the~\syn~16S rRNA anti-SD (\texttt{ACCUCCUUU}) and local mRNA secondary structure near each AUG.
For every predicted start site, the downstream reading frame was extended to the first in-frame stop codon (TAA or TAG) under genetic code~4, producing ${\sim}$26{,}000 candidate ORFs ($\geq$15 amino acids).

% Each ORF was classified by initiation mode: SD-dependent initiation ($\Delta G_\mathrm{mRNA{:}rRNA} < -7.2$~kcal/mol) was stratified into strong (TIR~$\geq$~10{,}000), moderate ($\geq$~1{,}000), or weak ($\geq$~100), while ORFs lacking significant SD complementarity were classified as SD-independent.

ORFs overlapping known sense reading frames were excluded, and the remainder were ranked by a synthesis rate score (isoform read count~$\times$~TIR).
The top 100 ORFs (of which ${\sim}$40\% exhibited strong SD-mediated initiation) were deduplicated by amino acid sequence and subjected to \textit{in silico} trypsin digestion (one missed cleavage; peptides of 7--25 residues).
Comparison of the resulting tryptic peptides against the canonical~\syn~proteome identified 44 novel ORFs producing at least one proteotypic peptide suitable for detection by bottom-up mass spectrometry.
These 44 novel ORFs were appended to the~\syn~protein database, and the proteomics raw data were re-searched with Spectronaut.